无线电测向把示向度误差限制在正负1度。本文将每条读数解释为半角1度的凸楔形，定位区域取诸楔之交；第二检测点沿第一示向度的正交基线布置；全向源采用9点侦听网和严格覆盖剪枝完成发现，再以距离上界收缩进入光学清除窗；定向源采用边长850 m的三角格网给出闭半平面瓣上的发现证书，听到后进入保扇形追逐。
问题一：合成正交交会例的对径圆覆盖为1；未与场地圆盘相交的两站有界样本覆盖率为1，相交裁剪后覆盖率为0.955；边长80 m的等边三角形对径圆覆盖为0。
问题二：第二检测点取轴上交角不小于45度的覆盖点，轴向距离800 m、横向偏移700 m。网格分位点横向偏移690.90909 m的事后直径中位数为48.2745 m，楔内70个样本上一次清除比例为0.2142857。固定轴向1100 m时，横向200 m的事后直径显著大于横向800 m。
问题三：9点网覆盖半径约956.051 m；有界收缩使每个已发现源最多再用7次动作清除，全任务动作上界为294。48个独立离线模型场景均全清；官方修订版测试尚待完成。
问题四：三角格网共37个侦听点，间距850 m，覆盖半径为490.7477288111819 m，小于500 m；平移量500 m时距离和为990.7477288111819 m。不合格间距上覆盖半径为519.6152422706632 m，不存在合格平移量。光学方格边长25 m的外接半径为17.677669529663685 m，小于20 m。
